\documentclass[11pt]{article}
\usepackage[margin=1in]{geometry}
\usepackage{booktabs}
\usepackage{hyperref}
\usepackage{longtable}
\usepackage{xcolor}
\usepackage{enumitem}
\usepackage{fancyvrb}
\hypersetup{colorlinks=true,linkcolor=blue,urlcolor=blue,citecolor=blue}

\title{Replication Report: Zhang, Pan \emph{et al.} (2023)\\
\emph{``Complete Genome Sequence and Pan-Genome Analysis of Shewanella oncorhynchi Z-P2, a Siderophore Putrebactin-Producing Bacterium''}}
\author{Ollie (OpenClaw AI) --- BVBRC Replication Project}
\date{2026-07-01}

\begin{document}
\maketitle

\begin{abstract}
This report documents an independent computational replication of the core comparative-genomics claims of Zhang, Pan \emph{et al.} (\emph{Microorganisms} 11(12):2961, 2023; DOI: 10.3390/microorganisms11122961; PMC10745600). The paper describes the first complete genome of \emph{Shewanella oncorhynchi} strain Z-P2 (GenBank CP132914 / RefSeq GCF\_030848765.1), a putrebactin-producing bacterium isolated from spoiled vacuum-packed crayfish. Using public NCBI data and a deliberately \emph{different} tool chain than the original (Prokka+Roary vs IPGA/PanOCT; fastANI vs MUMmer/kSNP-based ANI; RefSeq PGAP vs RAST), we reproduce genome size and GC content \emph{exactly}, tRNA/rRNA counts \emph{exactly}, closest-relative identity (YZ08), the sub-species ANI boundary conclusion, pan/core-genome cluster counts to within $\sim$1\% at a comparable orthology threshold, and all five secondary-metabolite BGCs (including the putrebactin NIS operon) at the paper's stated coordinates by their marker enzymes. Two claims (UPLC-MS m/z 373.21 for putrebactin; genomic-island/virulence/CRISPR re-annotation) were out of scope. Verdict: \textbf{PARTIAL REPLICATION (strong)}.
\end{abstract}

\section{Paper Summary}
Zhang \emph{et al.} report the first complete genome of the novel species \emph{Shewanella oncorhynchi}, strain Z-P2, isolated from spoiled vacuum-packed crayfish in Hubei, China. Z-P2 produces the hydroxamate siderophore \textbf{putrebactin} (a 20-membered macrocycle synthesized via the NRPS-independent siderophore (NIS) pathway). Sequencing used Nanopore PromethION + Illumina with Unicycler assembly. Annotation used RAST/COG/KEGG; BGC scanning used antiSMASH; comparative genomics used IPGA v1.09 / Prokka / PanOCT / kSNP / MUMmer against 10 \emph{S. putrefaciens} genomes. Putrebactin was confirmed by UPLC-MS.

\section{Claims Enumerated}
\begin{longtable}{@{}p{0.05\textwidth}p{0.55\textwidth}p{0.15\textwidth}p{0.15\textwidth}@{}}
\toprule
\# & Claim & Type & Tested here \\ \midrule
C1 & Single circular chromosome 5{,}034{,}612 bp; GC 45.4\% & Genome stats & Yes \\
C2 & 4544 CDS, 109 tRNA, 31 rRNA (11$\times$5S / 10$\times$16S / 10$\times$23S), 0 sRNA (RAST) & Annotation & Partial \\
C3 & 5 BGCs: APE, beta-lactone, putrebactin siderophore, EPA (hglE-KS), RiPP & BGC content & Yes (marker verification) \\
C4 & Pan-genome Z-P2 + 10 \emph{S. putrefaciens}: 9228 pan, 2681 core, 618 unique in Z-P2 & Comparative & Yes \\
C5 & Closest strain YZ08, ANI 90.09\%; all comparators $<$95\% species boundary & ANI/phylogeny & Yes \\
C6 & Putrebactin detected by UPLC-MS, [M+H]$^{+}$ m/z 373.21 & Wet-lab MS & \textbf{No (out of scope)} \\
C7 & 8 genomic islands, 30 virulence genes ($>$70\% id), 5 CRISPR arrays, 3 cas sets (subtype I-Fv) & Annotation detail & \textbf{No (not re-run)} \\
\bottomrule
\end{longtable}

\section{Method}
All work executed under \texttt{/data/stevens/bvbrc50/} on \texttt{uicgpu} (8$\times$A100, 255 cores) in conda env \texttt{bvbrc28}.
\begin{enumerate}[leftmargin=*]
  \item \textbf{Genome retrieval.} Strain resolved via \texttt{datasets summary genome taxon "Shewanella oncorhynchi"}; Z-P2 = GCF\_030848765.1. Downloaded Z-P2 + 10 complete \emph{S. putrefaciens} RefSeq genomes (including YZ08 = GCF\_019599085.1) via NCBI Datasets REST.
  \item \textbf{Genome statistics (C1, C2).} Length, contig count, GC recomputed from assembly FASTA (\texttt{work/genome\_stats.py}). tRNA/rRNA/CDS features counted from RefSeq PGAP GFF.
  \item \textbf{BGC verification (C3).} Extracted CDS products within each of the paper's five antiSMASH cluster coordinate windows from the Z-P2 RefSeq GFF; checked for canonical marker enzymes per cluster type. antiSMASH was \emph{not} re-run --- boundaries taken from the paper; content independently verified.
  \item \textbf{Pan/core-genome (C4).} Re-annotated all 11 genomes with Prokka 1.12 (\texttt{--genus Shewanella}); ran Roary 3.12 at BLASTp 95\% (default) and 70\% identity (closer to PanOCT-style orthology grouping). Cluster counts + per-genome uniques from \texttt{gene\_presence\_absence.csv}.
  \item \textbf{ANI/closest strain (C5).} fastANI, Z-P2 as query against all 11 genomes.
  \item \textbf{Scoring.} LLM judge (Argo \texttt{argo:gpt-5.2}, free) given full claim-by-claim evidence; verdict from the canonical vocabulary. No regex scoring.
\end{enumerate}
\textbf{Tool versions:} NCBI datasets, Prokka 1.12, Roary 3.12, fastANI, BLAST+, Python 3.11 --- all from conda env \texttt{bvbrc28}. \textbf{Endpoints:} free only (NCBI REST, Argo proxy).

\section{Results vs Paper}

\subsection{C1 --- Genome size \& GC (EXACT)}
\begin{tabular}{@{}lll@{}}
\toprule
Metric & Paper & Replication (GCF\_030848765.1) \\ \midrule
Chromosome length & 5{,}034{,}612 bp & 5{,}034{,}612 bp (exact) \\
GC content & 45.4\% & 45.40\% (exact) \\
Topology & single circular chromosome & 1 contig (consistent) \\
\bottomrule
\end{tabular}

\subsection{C2 --- Feature counts (tRNA/rRNA exact; CDS method-dependent)}
\begin{tabular}{@{}llll@{}}
\toprule
Feature & Paper (RAST) & Replication (RefSeq PGAP) & Match \\ \midrule
tRNA & 109 & 109 & exact \\
rRNA & 31 & 31 & exact \\
CDS & 4544 & 4290 (GFF) / 4220 (protein.faa) & $\sim$5.6\% lower \\
\bottomrule
\end{tabular}

tRNA/rRNA counts match \emph{exactly}. The CDS delta is the expected systematic gap between RAST and PGAP: RAST tends to call more/shorter ORFs and treats pseudogenes differently (RefSeq GFF also lists 42 pseudogenes). This is a pipeline artifact, not a genome discrepancy.

\subsection{C3 --- Five BGCs including putrebactin (all verified)}
Marker enzymes found in the RefSeq annotation at each paper-stated cluster window (evidence: \texttt{evidence/bgc\_regions.txt}):
\begin{longtable}{@{}p{0.20\textwidth}p{0.18\textwidth}p{0.50\textwidth}p{0.05\textwidth}@{}}
\toprule
BGC (paper) & Coordinates & Marker enzymes found & Verif. \\ \midrule
APE (aryl polyene) & 380{,}670--384{,}522 & $\beta$-ketoacyl-ACP synthase, 3-ketoacyl-ACP reductase FabG2, hotdog thioesterase & Yes \\
beta-lactone & 1{,}709{,}748--1{,}719{,}319 & acetyl/propionyl-CoA carboxylase ($\alpha$), HMG-CoA lyase, enoyl-CoA hydratase, carboxyl-transferase & Yes \\
Putrebactin (NIS) & 1{,}858{,}771--1{,}862{,}960 & IucA/IucC NIS synthetase, lysine N6-hydroxylase / L-ornithine N5-oxygenase, GNAT N-acetyltransferase, TonB-dependent siderophore receptor & Yes \\
EPA (hglE-KS) & 3{,}385{,}779--3{,}403{,}945 & eicosapentaenoate synthase PfaD, PfaB, phosphopantetheine-binding protein, hotdog thioesterase & Yes \\
RiPP & 4{,}156{,}609--4{,}163{,}174 & YcaO-domain protein (RiPP maturase), cupin, OsmC & Yes \\
\bottomrule
\end{longtable}

All five cluster \emph{types} independently confirmed at the paper's stated loci with the expected biosynthetic machinery. The putrebactin locus in particular carries the textbook NRPS-independent-siderophore enzyme set (ornithine monooxygenase $\rightarrow$ hydroxamate; IucA/IucC synthetase $\rightarrow$ macrocyclization; TonB receptor for ferri-putrebactin uptake).

\subsection{C4 --- Pan-genome (reproduced within $\sim$1\% at comparable threshold)}
\begin{tabular}{@{}llrrr@{}}
\toprule
Quantity & Paper (IPGA/PanOCT) & Roary @70\% id & $\Delta$ & Roary @95\% default \\ \midrule
Pan-genome clusters & 9228 & 9332 & +1.1\% & 17{,}326 \\
Core-genome clusters & 2681 & 2656 & --0.9\% & 684 \\
Z-P2 unique clusters & 618 & 531 & --14\% & 1756 \\
Unique-cluster range (all) & 24--640 & 0--774 & comparable & --- \\
\bottomrule
\end{tabular}

At a comparable orthology threshold (70\% id, closer to PanOCT grouping than Roary's strict 95\% default), pan- and core-genome cluster counts reproduce to within $\sim$1\% using a completely independent tool chain. Roary at 95\% over-splits near-orthologs, inflating pan-genome and singleton counts --- a controlled demonstration of the well-known sensitivity of pan-genome size to orthology threshold. The Z-P2 unique count (531 vs 618) is the same order of magnitude; the $\sim$14\% gap sits within expected cross-tool variance for singleton calling.

\subsection{C5 --- Closest strain \& ANI (closest-strain exact; ANI within $\sim$1.2 pts)}
\begin{tabular}{@{}rlr@{}}
\toprule
Rank & Strain & ANI to Z-P2 \\ \midrule
--- & Z-P2 (self) & 100.00\% \\
\textbf{1} & \textbf{YZ08 (GCF\_019599085.1)} & \textbf{91.25\%} \\
2 & FDAARGOS\_681 & 86.56\% \\
3 & WS13 & 86.48\% \\
\ldots & (six more) & 86.4--85.0\% \\
10 & NCTC12093 & 81.66\% \\
\bottomrule
\end{tabular}

Closest strain = YZ08 --- exact match to the paper. Paper ANI(Z-P2, YZ08) = 90.09\%; our fastANI value = 91.25\% --- within $\sim$1.2 percentage points, consistent with the different ANI algorithm (paper's MUMmer/kSNP-based ANI vs fastANI). All comparators fall below the 95\% species boundary, supporting the taxonomic conclusion that Z-P2 (\emph{S. oncorhynchi}) is a distinct species from \emph{S. putrefaciens}.

\subsection{C6, C7 --- Not tested}
\textbf{C6} (UPLC-MS m/z 373.21) is a wet-lab measurement, not computationally reproducible from public sequence data. \textbf{C7} (genomic islands / virulence / CRISPR counts) would require IslandPath-DIMOB, VFDB BLAST and CRISPRdigger re-runs and was not attempted in this pass. The RefSeq GFF does annotate 5 direct-repeat features and 12 riboswitches, broadly consistent with a CRISPR-bearing genome --- but this is \emph{not} a like-for-like reproduction of the paper's C7 numbers.

\section{LLM-Judge Assessment (Argo gpt-5.2, free)}
Per-claim ratings: C1 STRONG, C2 MODERATE, C3 MODERATE, C4 MODERATE, C5 STRONG --- \textbf{2 STRONG + 3 MODERATE, no FAIL, no contradictions}. Coverage: 5/7 claims tested ($\sim$71\% of computationally-testable claims). Verdict issued: \textbf{PARTIAL}. Full output: \texttt{evidence/llm\_judge\_argo\_gpt52.md}.

\section{Genuine Critique}
This section is deliberately adversarial. It records honestly what did \emph{not} replicate and where the replication itself is weak.

\subsection*{What did not replicate cleanly}
\begin{itemize}[leftmargin=*]
  \item \textbf{CDS count (C2) is off by 254 CDS ($\sim$5.6\%).} Paper: 4544 (RAST). Replication: 4290 (PGAP GFF) or 4220 (\texttt{protein.faa}). We attribute this to the RAST-vs-PGAP annotation-pipeline gap, but we did \emph{not} actually re-run RAST to confirm --- so this is an \emph{inferred} explanation, not a demonstrated one. A rigorous replication would re-run RAST on the same FASTA and show the count reconciling.
  \item \textbf{ANI value is off by $\sim$1.2 percentage points.} Paper: 90.09\% (MUMmer/kSNP-based). Replication: 91.25\% (fastANI). While the closest-strain identity is exact and the $<$95\% species-boundary conclusion is preserved, we did \emph{not} re-implement the paper's exact ANI algorithm. The gap is consistent with algorithmic differences, but this remains an approximate reproduction of the ANI value itself.
  \item \textbf{Z-P2 unique-cluster count (C4) is off by 14\%.} Paper: 618. Replication (Roary @70\%): 531. Singleton counts are notoriously sensitive to threshold choice, tool internals, and annotation-tool splits; we did not run IPGA/PanOCT directly to close this gap.
  \item \textbf{Two claims (C6, C7) are not tested at all.} C6 is genuinely wet-lab (UPLC-MS) and honestly out of scope for a purely computational replication. C7 is computationally reproducible in principle but was not attempted; the note that RefSeq annotates 5 direct-repeat features is \emph{not} equivalent evidence for 5 CRISPR arrays with 3 cas subtype-I-Fv sets --- it is a weak circumstantial signal, no more.
\end{itemize}

\subsection*{Where the replication design itself is weak}
\begin{itemize}[leftmargin=*]
  \item \textbf{BGC verification is coordinate-anchored to the paper.} antiSMASH was not re-run; we only verified that the paper's stated cluster \emph{windows} contain the expected marker enzymes. This is a check that the paper's cluster annotations are internally consistent, \emph{not} an independent BGC discovery. A stronger replication would re-run antiSMASH on the assembly and compare cluster boundaries and cluster-count independently.
  \item \textbf{Comparator set is inherited, not re-derived.} We used the same 10 \emph{S. putrefaciens} genomes referenced by the paper. If the paper's comparator choice biased pan-genome / ANI results, our replication inherits that bias.
  \item \textbf{No pseudogene reconciliation.} PGAP calls 42 pseudogenes on this genome; RAST typically merges some of these into the CDS pool. Without an explicit pseudogene bridge, the CDS-count gap remains partly unexplained.
  \item \textbf{LLM judge is a self-report, not a ground truth.} Per-claim ratings by \texttt{argo:gpt-5.2} are useful as a sanity check on the narrative but do not constitute independent scientific validation --- they are computed from our own evidence bundle.
  \item \textbf{No taxonomic re-adjudication.} We accept the paper's \emph{S. oncorhynchi} species designation as given. A stricter replication would compute 16S rRNA identity and dDDH against multiple \emph{Shewanella} type strains (\emph{S. baltica}, \emph{S. putrefaciens}, \emph{S. oncorhynchi} type) to independently support or contest the novel-species claim.
  \item \textbf{Coverage 5/7 = 71\% of \emph{computationally testable} claims, but only 5/7 = 71\% of \emph{all} claims.} The verdict PARTIAL is honest, but the untested claims (C7 especially) include ones we \emph{could} have addressed and did not.
\end{itemize}

\subsection*{What would strengthen this to FULL}
\begin{enumerate}[leftmargin=*]
  \item Re-run antiSMASH on GCF\_030848765.1; compare cluster count and boundaries.
  \item Re-run RAST on the same FASTA; reconcile the CDS-count gap explicitly.
  \item Run IslandPath-DIMOB, VFDB BLAST, and CRISPRCasFinder to test C7 like-for-like.
  \item Compute 16S rRNA identity and dDDH vs \emph{Shewanella} type strains to independently support the species boundary.
  \item Add a MUMmer-based ANI as a second ANI algorithm to bracket the fastANI value.
\end{enumerate}

\section{Conclusion}
Every computationally-testable core claim we reached was independently reproduced on real public data --- several \emph{exactly} (genome size, GC, tRNA, rRNA, closest-strain identity) and the rest \emph{quantitatively close} using deliberately different tools (fastANI vs kSNP/MUMmer; Prokka+Roary vs IPGA/PanOCT; RefSeq PGAP vs RAST). No claim was contradicted. The two untested claims are a wet-lab measurement (C6) and an annotation-detail re-run (C7), not the paper's central conclusions.

\vspace{0.5em}
\noindent\textbf{Verdict: PARTIAL REPLICATION (strong).}

\vspace{1em}
\noindent\emph{Artifacts:} \texttt{report/evidence/} (genome stats, fastANI table, Roary pan-genome analyses at 70\% and 95\%, BGC region gene content, LLM-judge output). \emph{Code + data pointers:} \texttt{work/}. Raw genomes, Prokka annotations, Roary outputs retained on uicgpu at \texttt{/data/stevens/bvbrc50/}.

\end{document}
